%% Question bank chapter file - generated by bank_tools.py build (do not edit while the bank is still being built:
%% it is regenerated from the per-test data in Question Bank\data). Every question is one bankq block:
%%   \begin{bankq}[id=..., source={...}, paper=.., qno=.., type=..., subtopic={...}, answer={...}, answer_src=key|solved,
%%                 working={...}, difficulty=1-10, diff_parts={K S I C F}, diff_note={...},
%%                 group=<passage id or empty>, also_in={duplicate copies in other tests}]
%% difficulty = min(10, K+S+I+C+F-1): K concepts 1-3, S steps 1-4, I insight 0-2, C computation 0-2, F format 0-1
%%   ...question LaTeX, options, figures (\bfig / \bfigi -> img/)...
%%   \end{bankq}
%% Passages shared by several questions are bankpassage blocks; their questions carry group=<passage id>.
\documentclass[10pt]{article}
\usepackage[a4paper,margin=1.8cm]{geometry}
\usepackage[T1]{fontenc}
\usepackage{amsmath,amssymb,mathtools}
\usepackage{graphicx}
\usepackage[export]{adjustbox}
\usepackage{xcolor,needspace,keyval}
\usepackage{enumitem}
\usepackage{booktabs,array,tabularx}
\usepackage[hidelinks]{hyperref}
\setlength{\parindent}{0pt}
\setlength{\parskip}{3pt}
\makeatletter
\define@key{bq}{id}{\def\bq@id{#1}}\define@key{bq}{source}{\def\bq@src{#1}}\define@key{bq}{paper}{\def\bq@paper{#1}}
\define@key{bq}{qno}{\def\bq@qno{#1}}\define@key{bq}{type}{\def\bq@type{#1}}\define@key{bq}{subtopic}{\def\bq@sub{#1}}
\define@key{bq}{answer}{\def\bq@ans{#1}}\define@key{bq}{answer_src}{\def\bq@asrc{#1}}\define@key{bq}{working}{\def\bq@work{#1}}
\define@key{bq}{group}{\def\bq@group{#1}}\define@key{bq}{also_in}{\def\bq@also{#1}}
\define@key{bq}{difficulty}{\def\bq@diff{#1}}\define@key{bq}{diff_parts}{\def\bq@dparts{#1}}\define@key{bq}{diff_note}{\def\bq@dnote{#1}}
\newenvironment{bankq}[1][]{\def\bq@work{}\def\bq@also{}\def\bq@group{}\def\bq@diff{?}\def\bq@dparts{}\def\bq@dnote{}\setkeys{bq}{#1}%
  \par\needspace{4\baselineskip}\medskip\noindent{\small\textbf{\bq@id}\hfill\color{gray}[\bq@src, Paper \bq@paper, Q\bq@qno]
  \bq@type\ $\cdot$ \bq@sub\ $\cdot$ \textbf{D\bq@diff}}\par\nopagebreak\relax}
  {\par{\small\color{teal}Answer: \textbf{\bq@ans}\ifx\bq@work\@empty\else\ (solved: \bq@work)\fi
  \ifx\bq@also\@empty\else\par\color{gray}Also in: \bq@also\fi}\par}
\define@key{bp}{id}{\def\bp@id{#1}}\define@key{bp}{source}{\def\bp@src{#1}}\define@key{bp}{questions}{\def\bp@qs{#1}}
\newenvironment{bankpassage}[1][]{\setkeys{bp}{#1}\par\needspace{4\baselineskip}\bigskip\noindent
  {\small\textbf{Passage \bp@id}\hfill\color{gray}[\bp@src; questions \bp@qs]}\par\nopagebreak\relax\itshape}{\par}
\makeatother
\newcommand{\bfig}[2]{\par\smallskip{\centering\includegraphics[width=#1pt,max width=0.62\linewidth,max height=0.32\textheight]{img/#2}\par}\smallskip}
\newcommand{\bfigi}[2]{\includegraphics[width=#1pt,max width=\linewidth,max height=0.32\textheight,valign=c]{img/#2}}
\newcommand{\bmissing}[1]{\textit{[figure missing in the source conversion: \detokenize{#1}]}}
\newenvironment{qoptions}{\begin{itemize}[leftmargin=1.8em,itemsep=1pt,topsep=2pt,label={}]}{\end{itemize}}
\newcommand{\opt}[2]{\item[(#1)] #2}
\providecommand{\fwstart}{}\providecommand{\fwend}{}\providecommand{\qfullsplit}{}
\newcolumntype{L}[1]{>{\raggedright\arraybackslash}p{#1}}
\newcolumntype{C}[1]{>{\centering\arraybackslash}p{#1}}
\newcolumntype{R}[1]{>{\raggedleft\arraybackslash}p{#1}}
\providecommand{\sectionhdr}[1]{}\providecommand{\subsectionhdr}[1]{}


\begin{document}
{\Large\bfseries Determinants --- Question Bank}\par\smallskip
{\small Mathematics chapter M08. Questions from all converted test papers, grouped by source series. Answers: printed key where it matches the paper, otherwise solved (answer\_src=solved, with working). Exact repeats across papers appear once, with also\_in listing the other copies.}\par\medskip\hrule

\section*{Score-I papers}
\begin{bankq}[id=SI02-P1-MATH-S3-Q03, source={Score-I - Test 02 (31 Dec 2021)}, paper=1, qno=3, type=single, subtopic={Properties}, answer={10}, answer_src=key, difficulty=7, diff_parts={K3 S1 I2 C2 F0}, diff_note={concepts: standard problem in this chapter, properties, apply the relevant results | steps: apply the relevant results; compute}]
If $\Delta=\begin{vmatrix}1&3\cos\theta&1\\\sin\theta&1&3\cos\theta\\1&\sin\theta&1\end{vmatrix}$, then maximum value of $\Delta$ is-
\end{bankq}
\begin{bankq}[id=SI02-P2-MATH-S1-Q03, source={Score-I - Test 02 (31 Dec 2021)}, paper=2, qno=3, type=single, subtopic={Properties}, answer={D}, answer_src=key, difficulty=4, diff_parts={K2 S1 I1 C1 F0}, diff_note={concepts: standard problem in this chapter, properties | steps: apply the relevant results; compute}]
If $a,b,c$ are unequal natural numbers, then the value of the determinant
\[\begin{vmatrix}a^2+b^2&ab+bc&ac+ba\\ba+cb&b^2+c^2&bc+ca\\ca+ab&cb+ac&c^2+a^2\end{vmatrix}\]
is
\begin{qoptions}
\opt{A}{$a+b+c$}
\opt{B}{$a+b$}
\opt{C}{$b+c$}
\opt{D}{$0$}
\end{qoptions}
\end{bankq}
\begin{bankq}[id=SI02-P2-MATH-S1-Q08, source={Score-I - Test 02 (31 Dec 2021)}, paper=2, qno=8, type=single, subtopic={Properties}, answer={A}, answer_src=key, difficulty=4, diff_parts={K2 S1 I1 C1 F0}, diff_note={concepts: standard problem in this chapter, properties | steps: apply the relevant results; compute}]
If $\begin{vmatrix}a&5x&p\\b&10y&5\\c&15z&15\end{vmatrix}=125$, then $\begin{vmatrix}3a&3b&c\\x&2y&z\\p&5&5\end{vmatrix}$ is equal to
\begin{qoptions}
\opt{A}{$25$}
\opt{B}{$50$}
\opt{C}{$75$}
\opt{D}{$100$}
\end{qoptions}
\end{bankq}
\begin{bankq}[id=SI02-P2-MATH-S2-Q02, source={Score-I - Test 02 (31 Dec 2021)}, paper=2, qno=2, type=numerical, subtopic={Properties}, answer={9}, answer_src=key, difficulty=8, diff_parts={K3 S1 I2 C2 F1}, diff_note={concepts: standard problem in this chapter, properties, apply the relevant results | steps: apply the relevant results; compute}]
If $\begin{vmatrix}1&x&x^2\\x&x^2&1\\x^2&1&x\end{vmatrix}=3$, then value of $\begin{vmatrix}x^3-1&0&x-x^4\\0&x-x^4&x^3-1\\x-x^4&x^3-1&0\end{vmatrix}$ is
\end{bankq}
\begin{bankq}[id=SI02-P2-MATH-S2-Q09, source={Score-I - Test 02 (31 Dec 2021)}, paper=2, qno=9, type=numerical, subtopic={Properties}, answer={0}, answer_src=key, difficulty=8, diff_parts={K3 S1 I2 C2 F1}, diff_note={concepts: standard problem in this chapter, properties, apply the relevant results | steps: apply the relevant results; compute}]
Let $a_1,a_2,a_3,\ldots$ be in G.P. Then the value of $\begin{vmatrix}\log a_n&\log a_{n+1}&\log a_{n+2}\\\log a_{n+3}&\log a_{n+4}&\log a_{n+5}\\\log a_{n+6}&\log a_{n+7}&\log a_{n+8}\end{vmatrix}$ is
\end{bankq}
\begin{bankq}[id=SI09-P1-MATH-S2-Q05, source={Score-I - Test 09 (16 Feb 2022)}, paper=1, qno=5, type=multi, subtopic={Properties}, answer={A,B,D}, answer_src=key, difficulty=8, diff_parts={K3 S1 I2 C2 F1}, diff_note={concepts: standard problem in this chapter, properties, apply the relevant results | steps: apply the relevant results; compute}]
For any $3\times3$ matrix $M$, let $|M|$ denote the determinant of $M$. Let
\begin{gather*}
E=\begin{bmatrix}1&2&3\\2&3&4\\8&13&18\end{bmatrix},\quad
P=\begin{bmatrix}1&0&0\\0&0&1\\0&1&0\end{bmatrix}\\
\text{and}\quad F=\begin{bmatrix}1&3&2\\8&18&13\\2&4&3\end{bmatrix}.
\end{gather*}
If $Q$ is a nonsingular matrix of order $3\times3$, then which of the following statements is (are) \textbf{TRUE} ?
\begin{qoptions}
\opt{A}{$F=PEP$ and $P^2=\begin{bmatrix}1&0&0\\0&1&0\\0&0&1\end{bmatrix}$}
\opt{B}{$|EQ+PFQ^{-1}|=|EQ|+|PFQ^{-1}|$}
\opt{C}{$\left|(EF)^3\right|>|EF|^2$}
\opt{D}{Sum of the diagonal entries of $P^{-1}EP+F$ is equal to the sum of diagonal entries of $E+P^{-1}FP$}
\end{qoptions}
\end{bankq}
\begin{bankq}[id=SI12-P1-MATH-S2-Q02, source={Score-I - Test 12 (09 Mar 2022)}, paper=1, qno=2, type=numerical, subtopic={Properties}, answer={1.50}, answer_src=key, difficulty=5, diff_parts={K1 S2 I1 C1 F1}, diff_note={concepts: determinant properties | steps: expand the determinant; it simplifies to a constant; so the sum of 25 terms is 150 | I: f is constant}]
If $f(x)=\begin{vmatrix}\cos(x+\alpha)&\cos(x+\beta)&\cos(x+\gamma)\\\sin(x+\alpha)&\sin(x+\beta)&\sin(x+\gamma)\\\sin(\beta-\gamma)&\sin(\gamma-\alpha)&\sin(\alpha-\beta)\end{vmatrix}$ and $f(2)=6$, if $k=\displaystyle\sum_{r=1}^{25}f(r)$, then find $\dfrac k{100}$
\end{bankq}

\section*{Score-II papers}
\begin{bankq}[id=SII03-P1-MATH-S1-Q06, source={Score-II - Test 03 (05 May 2021)}, paper=1, qno=6, type=single, subtopic={Properties}, answer={D}, answer_src=key, difficulty=4, diff_parts={K2 S1 I1 C1 F0}, diff_note={concepts: standard problem in this chapter, properties | steps: apply the relevant results; compute}]
$\begin{vmatrix}3x^2+19&2x-37&2\\x^2-29&x+31&3\\2x^2+17&3x&5\end{vmatrix}=ax^3+bx^2+cx+d$, then value of a is
\begin{qoptions}
\opt{A}{2}
\opt{B}{5}
\opt{C}{$-5$}
\opt{D}{$-8$}
\end{qoptions}
\end{bankq}
\begin{bankq}[id=SII04-P1-MATH-S3-Q05, source={Score-II - Test 04 (07 May 2021)}, paper=1, qno=5, type=numerical, subtopic={Cramer's rule}, answer={4}, answer_src=key, difficulty=8, diff_parts={K3 S1 I2 C2 F1}, diff_note={concepts: standard problem in this chapter, cramer's rule, apply the relevant results | steps: apply the relevant results; compute}]
If system of linear equations $(a-1)x+z=\alpha$, $x+(b-1)y=\beta$ and $y+(c-1)z=\lambda$ (where a, b, c $\in\mathbb{Z}$) does not have a unique solution, then maximum possible value of $|a+b+c|$ is:
\end{bankq}
\begin{bankq}[id=SII05-P1-MATH-S1-Q11, source={Score-II - Test 05 (22 Mar 2022)}, paper=1, qno=11, type=single, subtopic={Properties}, answer={C}, answer_src=key, difficulty=4, diff_parts={K2 S1 I1 C1 F0}, diff_note={concepts: standard problem in this chapter, properties | steps: apply the relevant results; compute}]
The value of\\ \resizebox{\linewidth}{!}{$\sin\frac\pi{14}\sin\frac{3\pi}{14}\sin\frac{5\pi}{14}\sin\frac{9\pi}{14}\sin\frac{11\pi}{14}\sin\frac{13\pi}{14}$}\\ is :
\begin{qoptions}
\opt{A}{$\dfrac1{16}$}
\opt{B}{$\dfrac1{64}$}
\opt{C}{$\dfrac1{128}$}
\opt{D}{None of these}
\end{qoptions}
\end{bankq}
\begin{bankq}[id=SII06-P1-MATH-S2-Q10, source={Score-II - Test 06 (05 Jul 2022)}, paper=1, qno=10, type=multi, subtopic={Properties}, answer={A,B,D}, answer_src=key, difficulty=8, diff_parts={K3 S1 I2 C2 F1}, diff_note={concepts: standard problem in this chapter, properties, apply the relevant results | steps: apply the relevant results; compute}]
Let $\lambda, \alpha, \beta \in \mathbb{R}$, consider the system of linear equations
$$10x + \lambda y = \alpha$$
$$-5x + 3y = 2\beta$$
then which of following statements is(are) correct?
\begin{qoptions}
\opt{A}{if $\lambda \neq -6$ then system has unique solutions for all values of $\alpha$ and $\beta$}
\opt{B}{if $\alpha + 4\beta = 0$ then system has infinite many solutions for $\lambda = -6$}
\opt{C}{if $\lambda = -6$ then system has infinite many solutions of all values of $\alpha$ and $\beta$}
\opt{D}{if $\alpha + 4\beta \neq 0$ then system has no solutions for $\lambda = -6$}
\end{qoptions}
\end{bankq}
\begin{bankq}[id=SII08-P1-MATH-S1-Q03, source={Score-II - Test 08 (05 Feb 2023)}, paper=1, qno=3, type=multi, subtopic={Properties}, answer={B,D}, answer_src=key, difficulty=8, diff_parts={K3 S1 I2 C2 F1}, diff_note={concepts: standard problem in this chapter, properties, apply the relevant results | steps: apply the relevant results; compute}]
Let $|M|$ denote the determinant of a square matrix $M$. Let $g:\left[0,\dfrac{\pi}{2}\right]\to\mathbb{R}$ be the function defined by
$$g(\theta) = \sqrt{f(\theta)-1} + \sqrt{f\left(\dfrac{\pi}{2}-\theta\right)-1}\ \text{ where}$$
\par\noindent\resizebox{\linewidth}{!}{$f(\theta) = \dfrac12\begin{vmatrix}1 & \sin\theta & 1\\ -\sin\theta & 1 & \sin\theta\\ -1 & -\sin\theta & 1\end{vmatrix} + \begin{vmatrix}\sin\pi & \cos\left(\theta+\dfrac{\pi}{4}\right) & \tan\left(\theta-\dfrac{\pi}{4}\right)\\ \sin\left(\theta-\dfrac{\pi}{4}\right) & -\cos\dfrac{\pi}{2} & \log_e\left(\dfrac{4}{\pi}\right)\\ \cot\left(\theta+\dfrac{\pi}{4}\right) & \log_e\left(\dfrac{4}{\pi}\right) & \tan\pi\end{vmatrix}.$}\par
Let $p(x)$ be a quadratic polynomial whose roots are the maximum and minimum values of the function $g(\theta)$, and $p(2) = 2-\sqrt2$. Then, which of the following is/are TRUE?
\begin{qoptions}
\opt{A}{$p\left(\dfrac{5-\sqrt2}{4}\right)<0$}
\opt{B}{$p\left(\dfrac{5\sqrt2-1}{4}\right)>0$}
\opt{C}{$p\left(\dfrac{1+3\sqrt2}{4}\right)>0$}
\opt{D}{$p\left(\dfrac{3+\sqrt2}{4}\right)<0$}
\end{qoptions}
\end{bankq}
\begin{bankq}[id=SII08-P2-MATH-S3-Q04, source={Score-II - Test 08 (05 Feb 2023)}, paper=2, qno=4, type=numerical, subtopic={Properties}, answer={4}, answer_src=key, difficulty=8, diff_parts={K3 S1 I2 C2 F1}, diff_note={concepts: standard problem in this chapter, properties, apply the relevant results | steps: apply the relevant results; compute}]
Let $\beta$ be a real number. Consider the matrix $A = \begin{pmatrix}\beta & 0 & 1\\ 2 & 1 & -2\\ 3 & 1 & -2\end{pmatrix}$. If $A^7-(\beta-1)A^6-\beta A^5$ is a singular matrix, then the value of $3(\beta+1)$ is\_\_\_\_\_\_\_.
\end{bankq}
\begin{bankq}[id=SII09-P1-MATH-S1-Q06, source={Score-II - Test 09 (12 Feb 2023)}, paper=1, qno=6, type=multi, subtopic={Properties}, answer={A,B,C}, answer_src=key, difficulty=8, diff_parts={K3 S1 I2 C2 F1}, diff_note={concepts: standard problem in this chapter, properties, apply the relevant results | steps: apply the relevant results; compute}]
The sum of absolute values of x satisfying equation
$$\begin{vmatrix}2^{11} & 2^{11} & -2^{12}-x\\ -2^{11}+x & 2^{12} & -2^{11}\\ 2^{12} & -2^{11}+x & -2^{11}\end{vmatrix}=0$$
is $\alpha$, then
\begin{qoptions}
\opt{A}{$\alpha$ is divisible by 3}
\opt{B}{$\alpha$ is divisible by 8}
\opt{C}{Number of divisors of $\alpha$ = 26}
\opt{D}{$\alpha$ is divisible by 54}
\end{qoptions}
\end{bankq}
\begin{bankq}[id=SII13-P1-MATH-S1-Q12, source={Score-II - Test 13 (06 Feb 2024)}, paper=1, qno=12, type=single, subtopic={Cramer's rule}, answer={B}, answer_src=key, difficulty=4, diff_parts={K2 S1 I1 C1 F0}, diff_note={concepts: standard problem in this chapter, cramer's rule | steps: apply the relevant results; compute}]
The set of all values of $\lambda$ for which the system of linear equations.\\
$x-2y-2z=\lambda x$\\
$x+2y+z=\lambda y$\\
$-x-y=\lambda z$\\
has a non-trivial solution.
\begin{qoptions}
\opt{A}{contains more than two elements}
\opt{B}{is a singleton}
\opt{C}{is an empty set}
\opt{D}{contains exactly two elements}
\end{qoptions}
\end{bankq}
\begin{bankq}[id=SII15-P1-MATH-S1-Q01, source={Score-II - Test 15 (09 Feb 2024)}, paper=1, qno=1, type=single, subtopic={Properties}, answer={A}, answer_src=key, difficulty=4, diff_parts={K2 S1 I1 C1 F0}, diff_note={concepts: standard problem in this chapter, properties | steps: apply the relevant results; compute}]
If $A$ is square matrix of order 3 such that $|A|=3$, then $|\text{adj}\,(\text{adj}\,(\text{adj}\,A))|$ is-
\begin{qoptions}
\opt{A}{$3^8$}
\opt{B}{$3^9$}
\opt{C}{$3^{10}$}
\opt{D}{$3^{11}$}
\end{qoptions}
\end{bankq}
\begin{bankq}[id=SII15-P1-MATH-S1-Q08, source={Score-II - Test 15 (09 Feb 2024)}, paper=1, qno=8, type=single, subtopic={Properties}, answer={D}, answer_src=key, difficulty=4, diff_parts={K2 S1 I1 C1 F0}, diff_note={concepts: standard problem in this chapter, properties | steps: apply the relevant results; compute}]
If $\begin{vmatrix}1&x&x^2\\x&x^2&1\\x^2&1&x\end{vmatrix}=3$ then the value of $\begin{vmatrix}x^3-1&0&x-x^4\\0&x-x^4&x^3-1\\x-x^4&x^3-1&0\end{vmatrix}$ is
\begin{qoptions}
\opt{A}{0}
\opt{B}{3}
\opt{C}{6}
\opt{D}{9}
\end{qoptions}
\end{bankq}
\begin{bankq}[id=SII15-P1-MATH-S1-Q13, source={Score-II - Test 15 (09 Feb 2024)}, paper=1, qno=13, type=single, subtopic={Properties}, answer={D}, answer_src=key, difficulty=4, diff_parts={K2 S1 I1 C1 F0}, diff_note={concepts: standard problem in this chapter, properties | steps: apply the relevant results; compute}]
If $a,b,c$ are real then the value of determinant $\begin{vmatrix}a^2+1&ab&ac\\ab&b^2+1&bc\\ac&bc&c^2+1\end{vmatrix}=1$ if
\begin{qoptions}
\opt{A}{$a+b+c=2$}
\opt{B}{$a+b+c=1$}
\opt{C}{$a+b+c=-1$}
\opt{D}{$a=b=c=0$}
\end{qoptions}
\end{bankq}
\begin{bankq}[id=SII15-P2-MATH-S2-Q10, source={Score-II - Test 15 (09 Feb 2024)}, paper=2, qno=10, type=multi, subtopic={Properties}, answer={A,B,C}, answer_src=key, difficulty=8, diff_parts={K3 S1 I2 C2 F1}, diff_note={concepts: standard problem in this chapter, properties, apply the relevant results | steps: apply the relevant results; compute}]
Let $\Delta=\begin{vmatrix}a&b&c\\a-b&b-c&c-a\\b+c&c+a&a+b\end{vmatrix}$.\\
Which option is/are always correct
\begin{qoptions}
\opt{A}{If $a,b,c$ are unequal positive real numbers, then $\Delta$ is always positive}
\opt{B}{If $\Delta=0$, then $a^3+b^3+c^3=3abc$}
\opt{C}{If $a,b,c$ are unequal negative real numbers, then $\Delta$ is always negative}
\opt{D}{If $\Delta=0$, then $a^2+b^2+c^2=ab+bc+ca$}
\end{qoptions}
\end{bankq}
\begin{bankq}[id=SII17-P1-MATH-S2-Q02, source={Score-II - Test 17 (13 Feb 2024)}, paper=1, qno=2, type=numerical, subtopic={Properties}, answer={9}, answer_src=solved, working={|3A\^{}-1|=27/9=3; |adj M|=|M|\^{}2=9.}, difficulty=8, diff_parts={K3 S1 I2 C2 F1}, diff_note={concepts: standard problem in this chapter, properties, apply the relevant results | steps: apply the relevant results; compute}]
If $|A|_{3\times3}=9$, then value of $\left|\text{adj}\left(\dfrac A3\right)^{-1}\right|$ is $-$
\end{bankq}
\begin{bankq}[id=SII18-P1-MATH-S2-Q08, source={Score-II - Test 18 (13 Feb 2024)}, paper=1, qno=8, type=multi, subtopic={Properties}, answer={B,C,D}, answer_src=key, difficulty=8, diff_parts={K3 S1 I2 C2 F1}, diff_note={concepts: standard problem in this chapter, properties, apply the relevant results | steps: apply the relevant results; compute}]
For any $3\times3$ matrix M, let $|$M$|$ denote the determinant of M. Let A = $\begin{bmatrix}1&3&9\\2&2&10\\3&1&11\end{bmatrix}$, B = $\begin{bmatrix}10&2&2\\11&3&1\\9&1&3\end{bmatrix}$, C = $\begin{bmatrix}0&1&0\\0&0&1\\1&0&0\end{bmatrix}$. If D = $\begin{bmatrix}1&5&2\\3&7&4\\2&m&1\end{bmatrix}$ is a matrix such that $m\geqslant2$, then which of the following statements is(are) CORRECT?

\begin{qoptions}
\opt{A}{$|(AD)^{2024}| > |AD|^{2023}$}
\opt{B}{$|AD+C^2BD^{-1}| = |AD| + |CBD^{-1}|$}
\opt{C}{B = CAC and $C^3 = \begin{bmatrix}1&0&0\\0&1&0\\0&0&1\end{bmatrix}$}
\opt{D}{$C^{2024}B^2C^{2024} = AC^2A$}
\end{qoptions}
\end{bankq}
\begin{bankq}[id=SII21-P1-MATH-S1-Q06, source={Score-II - Test 21 (14 Mar 2024)}, paper=1, qno=6, type=multi, subtopic={Properties}, answer={B,D}, answer_src=key, difficulty=8, diff_parts={K3 S1 I2 C2 F1}, diff_note={concepts: standard problem in this chapter, properties, apply the relevant results | steps: apply the relevant results; compute}]
Let $|M|$ denote the determinant of a square matrix $M$. Let g$:\left[0,\dfrac{\pi}{2}\right]\to\mathbb R$ be the function defined by
\[g(\theta) = \sqrt{f(\theta)-1} + \sqrt{f\left(\dfrac{\pi}{2}-\theta\right)-1}\ \text{where}\]
\par\noindent\resizebox{\linewidth}{!}{$\displaystyle f(\theta) = \dfrac12\begin{vmatrix}1 & \sin\theta & 1\\ -\sin\theta & 1 & \sin\theta\\ -1 & -\sin\theta & 1\end{vmatrix} + \begin{vmatrix}\sin\pi & \cos\left(\theta+\dfrac{\pi}{4}\right) & \tan\left(\theta-\dfrac{\pi}{4}\right)\\ \sin\left(\theta-\dfrac{\pi}{4}\right) & -\cos\dfrac{\pi}{2} & \log_e\left(\dfrac{4}{\pi}\right)\\ \cot\left(\theta+\dfrac{\pi}{4}\right) & \log_e\left(\dfrac{\pi}{4}\right) & \tan\pi\end{vmatrix}$}\par
Let p($x$) be a quadratic polynomial whose roots are the maximum and minimum values of the function g($\theta$), and p(2) $=4-2\sqrt2$. Then, which of the following is/are TRUE ?
\begin{qoptions}
\opt{A}{p$\left(\dfrac{3+\sqrt2}{4}\right)>0$}
\opt{B}{p$\left(\dfrac{1+3\sqrt2}{4}\right)<0$}
\opt{C}{p$\left(\dfrac{5\sqrt2-1}{4}\right)<0$}
\opt{D}{p$\left(\dfrac{5-\sqrt2}{4}\right)>0$}
\end{qoptions}
\end{bankq}
\begin{bankq}[id=SII22-P1-MATH-S1-Q03, source={Score-II - Test 22 (14 Mar 2024)}, paper=1, qno=3, type=single, subtopic={Properties}, answer={B}, answer_src=key, difficulty=4, diff_parts={K2 S1 I1 C1 F0}, diff_note={concepts: standard problem in this chapter, properties | steps: apply the relevant results; compute}]
If A is matrix of order $(3\times3)$ and $|A|=2$ then value of $|adj\ adj\ adj\ adj\ adj\ (2A)|$ is equal to : --
\begin{qoptions}
\opt{A}{$2^9$}
\opt{B}{$2^{128}$}
\opt{C}{$2^{32}$}
\opt{D}{None}
\end{qoptions}
\end{bankq}
\begin{bankq}[id=SII22-P1-MATH-S2-Q01, source={Score-II - Test 22 (14 Mar 2024)}, paper=1, qno=1, type=numerical, subtopic={Properties}, answer={1}, answer_src=key, difficulty=8, diff_parts={K3 S1 I2 C2 F1}, diff_note={concepts: standard problem in this chapter, properties, apply the relevant results | steps: apply the relevant results; compute}]
The number of distinct real solution of\\
\resizebox{\linewidth}{!}{$\begin{vmatrix}\sin x&\cos x&\cos x\\\cos x&\sin x&\cos x\\\cos x&\cos x&\sin x\end{vmatrix}=0$}\\
in the interval $-\dfrac\pi4\le x\le\dfrac\pi4$ is-
\end{bankq}
\begin{bankq}[id=SII23-P1-MATH-S3-Q03, source={Score-II - Test 23 (17 Mar 2024)}, paper=1, qno=3, type=numerical, subtopic={Cramer's rule}, answer={2.665}, answer_src=key, difficulty=8, diff_parts={K3 S1 I2 C2 F1}, diff_note={concepts: standard problem in this chapter, cramer's rule, apply the relevant results | steps: apply the relevant results; compute}]
Let $(x_0,y_0,z_0)$ be the solution of equation
\begin{align*}
3x-4y+\dfrac{1}{xy}&=12\\
4z-12x+\dfrac{1}{zx}&=3\\
12y-3z+\dfrac{1}{yz}&=4
\end{align*}
Then $\dfrac{z_0}{4x_0y_0}$ is
\end{bankq}
\begin{bankq}[id=SII25-P1-MATH-S2-Q09, source={Score-II - Test 25 (20 Mar 2024)}, paper=1, qno=9, type=numerical, subtopic={Properties}, answer={2}, answer_src=key, difficulty=8, diff_parts={K3 S1 I2 C2 F1}, diff_note={concepts: standard problem in this chapter, properties, apply the relevant results | steps: apply the relevant results; compute}]
Entries of a $2\times2$ determinant are chosen from the set $\{-1,1\}$. The probability that determinant has zero value is equal to $P$ then $4P$ is \underline{\hspace{2em}}.
\end{bankq}
\begin{bankq}[id=SII26-P1-MATH-S2-Q05, source={Score-II - Test 26 (20 Mar 2024)}, paper=1, qno=5, type=multi, subtopic={Properties}, answer={B,D}, answer_src=key, difficulty=8, diff_parts={K3 S1 I2 C2 F1}, diff_note={concepts: standard problem in this chapter, properties, apply the relevant results | steps: apply the relevant results; compute}]
Let $n\times n$ matrix is A$_n$ whose all elements are n except diagonal elements and given that a$_{ii}$ = i, for i = 1, 2, 3, ... n; then
\begin{qoptions}
\opt{A}{$|A_n|=n|A_{n-1}|$}
\opt{B}{$|A_n|=-n|A_{n-1}|$}
\opt{C}{$|A_n|=(-1)^n\,n!$}
\opt{D}{$|A_n|=(-1)^{n-1}\,n!$}
\end{qoptions}
\end{bankq}
\begin{bankq}[id=SII35-P1-MATH-S3-Q06, source={Score-II - Test 35 (04 Feb 2025)}, paper=1, qno=6, type=numerical, subtopic={Properties}, answer={2}, answer_src=key, difficulty=8, diff_parts={K3 S1 I2 C2 F1}, diff_note={concepts: standard problem in this chapter, properties, apply the relevant results | steps: apply the relevant results; compute}]
Suppose $\det\begin{bmatrix}\displaystyle\sum_{k=1}^{n}k&\displaystyle\sum_{k=1}^{n}{}^nC_kk^2\\ \displaystyle\sum_{k=1}^{n}{}^nC_kk&\displaystyle\sum_{k=0}^{n}{}^nC_k3^k\end{bmatrix}+\dfrac{n(n+1).4^n}2=0$, holds for some positive integer n. Then $\dfrac9{251}\displaystyle\sum_{k=1}^{n}\dfrac{{}^nC_k}{k+1}$ equals
\end{bankq}
\begin{bankq}[id=SII44-P1-MATH-S1-Q02, source={Score-II - Test 44 (23 Mar 2025)}, paper=1, qno=2, type=single, subtopic={Properties}, answer={Bonus}, answer_src=solved, working={Printed key covers only one subject (chemistry) and this question was not solved; banked as Bonus.}, difficulty=7, diff_parts={K3 S1 I2 C2 F0}, diff_note={concepts: standard problem in this chapter, properties, apply the relevant results | steps: apply the relevant results; compute}]
If $\displaystyle\prod_{r=0}^{89}\left(a+b\omega^r+c\omega^{2r}\right)=\begin{vmatrix}a&b&c\\b&c&a\\c&a&b\end{vmatrix}^k$; then k is equal to (where $\omega$ is the complex (non real) cube root of unity and a,b,c$\in\mathbb R$)

\begin{qoptions}
\opt{A}{20}
\opt{B}{30}
\opt{C}{21}
\opt{D}{31}
\end{qoptions}
\end{bankq}
\begin{bankq}[id=SII47-P1-MATH-S1-Q19, source={Score-II - Test 47 (30 Mar 2025)}, paper=1, qno=19, type=single, subtopic={Cramer's rule}, answer={A}, answer_src=key, difficulty=4, diff_parts={K2 S1 I1 C1 F0}, diff_note={concepts: standard problem in this chapter, cramer's rule | steps: apply the relevant results; compute}]
Suppose $p,q,r\ne0$ and the system of equations
\[(p+a)x+by+cz=0\]
\[ax+(q+b)y+cz=0\]
\[ax+by+(r+c)z=0\]
has non-trivial solutions, then the value of $\dfrac{a}{p}+\dfrac{b}{q}+\dfrac{c}{r}$ is
\begin{qoptions}
\opt{A}{$-1$}
\opt{B}{0}
\opt{C}{1}
\opt{D}{2}
\end{qoptions}
\end{bankq}

\section*{Score Advanced papers}
\begin{bankq}[id=SA02-P1-MATH-S3-Q06, source={Score Advanced - Test 02 (13 May 2021)}, paper=1, qno=6, type=numerical, subtopic={Properties}, answer={1.4524}, answer_src=solved, working={Determinant simplifies to -3cos x/4+cos2x/2-cos3x/4-1/2; integral = -(2/3+pi/4)}, difficulty=7, diff_parts={K2 S2 I2 C1 F1}, diff_note={concepts: determinant integral, properties | steps: simplify the determinant; integrate; evaluate}]
Let $f(x) =$\par{\scriptsize$\begin{vmatrix} \sec x & \cos x & \sec^2 x + \cot x\,\mathrm{cosec}\,x + \cos x\\ \cos^2 x & \cos^2 x & \mathrm{cosec}^2 x + \cos^4 x\\ 1 & \cos^2 x & \cos^2 x \end{vmatrix}$}\par then the value $\left|\displaystyle\int_0^{\pi/2} f(x)\,dx\right|$ is (use $\pi = \dfrac{22}{7}$)
\end{bankq}
\begin{bankq}[id=SA03-P1-MATH-S1-Q02, source={Score Advanced - Test 03 (27 Jun 2021)}, paper=1, qno=2, type=multi, subtopic={Properties}, answer={B,C}, answer_src=solved, working={det=-d\^{}2(12a+13d)}, difficulty=5, diff_parts={K2 S1 I1 C1 F1}, diff_note={concepts: determinant with arithmetic terms, properties | steps: expand; simplify}]
Let A = $\begin{vmatrix} a & a + d & a + 3d\\ a + d & a + 2d & a\\ a + 2d & a & a + d \end{vmatrix}$, then
\begin{qoptions}
\opt{A}{A = 0}
\opt{B}{A depends on a}
\opt{C}{A depends on d}
\opt{D}{A is independent of a and d}
\end{qoptions}
\end{bankq}
\begin{bankq}[id=SA03-P1-MATH-S3-Q03, source={Score Advanced - Test 03 (27 Jun 2021)}, paper=1, qno=3, type=matching, subtopic={Properties}, answer={P$\to$1, Q$\to$3, R$\to$2}, answer_src=solved, working={sum of sin\^{}2=2; detC=16; lambda=4}, difficulty=8, diff_parts={K3 S2 I2 C1 F1}, diff_note={concepts: determinant and matrix power identities, properties, expand | steps: expand; determinant of products; power of a matrix}]
\noindent{\footnotesize\setlength{\tabcolsep}{2pt}\begin{tabular}{|L{.7cm}|L{4.0cm}|L{.7cm}|L{1.4cm}|}\hline
 & List-I & & List-II\\\hline
(P) & $\begin{vmatrix} 1 & \cos\alpha & \cos\beta\\ \cos\alpha & 1 & \cos\gamma\\ \cos\beta & \cos\gamma & 1 \end{vmatrix}$ = $\begin{vmatrix} 0 & \cos\alpha & \cos\beta\\ \cos\alpha & 0 & \cos\gamma\\ \cos\beta & \cos\gamma & 0 \end{vmatrix}$, then $\sin^2\alpha + \sin^2\beta + \sin^2\gamma$ = & (1) & 2\\\hline
(Q) & If A = $\begin{pmatrix} 3 & 1\\ -1 & 1 \end{pmatrix}$ and a matrix is defined as C = BAB$^{-1}$(B$^{-1}$A$^T$B), then detC is equal to & (2) & 4\\\hline
(R) & If A = $\begin{pmatrix} 1 & 1\\ -1 & 1 \end{pmatrix}$ and A$^4$ = $-\lambda$I, then $\lambda$ equals to & (3) & 16\\\hline
\end{tabular}}\par
\end{bankq}
\begin{bankq}[id=SA03-P2-MATH-S1-Q19, source={Score Advanced - Test 03 (27 Jun 2021)}, paper=2, qno=19, type=single, subtopic={Properties}, answer={B}, answer_src=solved, working={36 lambda\^{}6 gives a=6 and lambda=2}, difficulty=6, diff_parts={K2 S2 I2 C1 F0}, diff_note={concepts: determinant of the adjoint of a scaled matrix, properties | steps: det A=6; adjoint rule; compare}]
If $A=\begin{bmatrix}3&2&1\\1&2&1\\0&1&2\end{bmatrix}$ also, $|\mathrm{adj}\,\lambda A|=\lambda^{a}.a^{\lambda}$ (where $a,\lambda\in N$). Then the value of `$\alpha+\lambda$', is-
\begin{qoptions}
\opt{A}{6}
\opt{B}{8}
\opt{C}{10}
\opt{D}{12}
\end{qoptions}
\end{bankq}
\begin{bankq}[id=SA03-P2-MATH-S2-Q02, source={Score Advanced - Test 03 (27 Jun 2021)}, paper=2, qno=2, type=numerical, subtopic={Expansion}, answer={2}, answer_src=solved, working={det=3 for every x}, difficulty=6, diff_parts={K2 S1 I2 C1 F1}, diff_note={concepts: ratio from a determinant, expansion | steps: simplify the determinant; evaluate}]
If $f(x)=$ {\scriptsize$\begin{vmatrix}3x^2-6x+5&6x^2+6x-3&3x^2+3x-2\\2x-2&4x+2&2x+1\\1&2&1\end{vmatrix}$}, then the value of $\dfrac{f(1)+f(3)}{f(2)}$
\end{bankq}
\begin{bankq}[id=SA05-P1-MATH-S1-Q06, source={Score Advanced - Test 05 (15 Jul 2021)}, paper=1, qno=6, type=single, subtopic={Expansion}, answer={C}, answer_src=key, difficulty=6, diff_parts={K2 S2 I2 C1 F0}, diff_note={concepts: sum of a tangent over the roots of a determinant, expansion | steps: expand the determinant to cos\^{}3x-sin\^{}3x; tan x=1; evaluate}]
If x $\in$ [0, 2$\pi$] and $\begin{vmatrix} 0 & \cos x & -\sin x\\ \sin x & 0 & \cos x\\ \cos x & \sin x & 0\end{vmatrix} = 0$, then $\sum\tan\left(\dfrac{\pi}{3}+x\right)$ is equal to-
\begin{qoptions}
\opt{A}{$\sqrt{3}-2$}
\opt{B}{$4+2\sqrt{3}$}
\opt{C}{$-2(2+\sqrt{3})$}
\opt{D}{$4-2\sqrt{3}$}
\end{qoptions}
\end{bankq}
\begin{bankq}[id=SA07-P2-MATH-S3-Q07, source={Score Advanced - Test 07 (16 Apr 2023)}, paper=2, qno=7, type=single, subtopic={Properties}, answer={5}, answer_src=key, difficulty=7, diff_parts={K3 S1 I2 C2 F0}, diff_note={concepts: global minimum of a trace of an inverse adjoint, properties, det A=3 gives | steps: det A=3 gives}]
Let A = $\begin{bmatrix}\text{x}^3+1&\text{x}+5&3\text{x}+2\\ \text{y}+1&-6\text{x}^2+2&\text{z}-1\\ 2&3&9\text{x}+6\end{bmatrix}$ and |A| = 3\par
If $f$(x) = tr(B$^{-1}$) and B = adj(A), then global minimum value of $f$(x) in x $\in$[0, 6] is p and it is achieved at `n' number of value(s) of x then n + p is equal to
\end{bankq}
\begin{bankq}[id=SA08-P1-MATH-S3-Q06, source={Score Advanced - Test 08 (22 Apr 2023)}, paper=1, qno=6, type=single, subtopic={Properties}, answer={2}, answer_src=key, difficulty=5, diff_parts={K2 S1 I2 C1 F0}, diff_note={concepts: determinant relations of adjugates, properties | steps:}]
If A and B are non-singular matrices of order three such that adj(AB) = $\begin{bmatrix}1&1&1\\1&\text{p}&1\\1&1&\text{p}\end{bmatrix}$ and $|\text{B}^2\text{adj(A)}|$ = p$^2$ $-$ 3, then p is
\end{bankq}
\begin{bankq}[id=SA08-P2-MATH-S3-Q05, source={Score Advanced - Test 08 (22 Apr 2023)}, paper=2, qno=5, type=numerical, subtopic={Properties}, answer={0}, answer_src=key, difficulty=8, diff_parts={K3 S1 I2 C2 F1}, diff_note={concepts: determinant of a combination with the adjugate, properties, evaluate matrix entries | steps: evaluate matrix entries; det(2M\^{}T+adj M)}]
Let matrix M = \par\noindent\resizebox{\linewidth}{!}{$\begin{bmatrix}\tan\left(\frac{301\pi}{3}\right) & \sec(2016\pi) & \cot\left(\frac{2015\pi}{2}\right)\\ \cot\left(\tan^{-1}\frac{1}{2}+\tan^{-1}\frac{1}{3}\right) & 2\sin\left(\frac{4\pi}{3}\right) & \sec^{-1}(2016)-\cos^{-1}\left(\frac{1}{2016}\right)\\ \cos\left(\frac{2009\pi}{2}\right) & \cos^{-1}(1) & \sec\left(\frac{301\pi}{3}\right)\end{bmatrix}$}\par , then det.(2M$^\text{T}$ + adj.M) is equal to
\end{bankq}
\begin{bankq}[id=SA09-P1-MATH-S2-Q06, source={Score Advanced - Test 09 (29 Apr 2023)}, paper=1, qno=6, type=multi, subtopic={Expansion}, answer={A,C,D}, answer_src=key, difficulty=8, diff_parts={K3 S1 I2 C2 F1}, diff_note={concepts: determinant combination that makes the integral minimal, expansion, combine the three determinants into a quadratic in g | steps: combine the three determinants into a quadratic in g; minimise the integral}]
Consider f(x) = \par\noindent\resizebox{\linewidth}{!}{$\begin{vmatrix}5&7&3\\2&9g^2(\text{x})&2\text{x}^2\\\text{x}&3g(\text{x})-\text{x}^2&1\end{vmatrix}-\begin{vmatrix}2&4&5\\2\text{x}^2&2&9g^2(\text{x})\\1&\text{x}&3g(\text{x})-\text{x}^2\end{vmatrix}+\begin{vmatrix}2&\text{x}&0\\9g^2(\text{x})&3g(\text{x})-\text{x}^2&-2\\2\text{x}^2&1&-1\end{vmatrix}$}\par
(where g(x) is polynomial function) such that m = $\displaystyle\int_0^1 f(\text{x})\text{dx}$ is minimum. Then
\begin{qoptions}
\opt{A}{m = 1/5}
\opt{B}{m = 2/5}
\opt{C}{Number of point of inflection of $\phi(\text{x}) = \displaystyle\int_0^\text{x}\left[\left(\int_0^\text{t} g(\text{y})\text{dy}-\left(\text{e}^\text{t}-\text{e}^{-\text{t}}\right)\right)\right]\text{dt}$ is 0}
\opt{D}{g(x) $\leqslant$ cos x $\forall$ x $\in$ [$-$1, 1]}
\end{qoptions}
\end{bankq}
\begin{bankq}[id=SA09-P2-MATH-S1-Q07, source={Score Advanced - Test 09 (29 Apr 2023)}, paper=2, qno=7, type=multi, subtopic={Properties}, answer={B,C,D}, answer_src=key, difficulty=8, diff_parts={K3 S1 I2 C2 F1}, diff_note={concepts: sign of a structured determinant, properties, the determinant is a Vandermonde-type product | steps: the determinant is a Vandermonde-type product}]
Consider the determinant\par\noindent\resizebox{\linewidth}{!}{$\text{D}=\begin{vmatrix}\dfrac{1-\ell_1^3\text{m}_1^3}{1-\ell_1\text{m}_1}&\dfrac{1-\ell_2^3\text{m}_1^3}{1-\ell_2\text{m}_1}&\dfrac{1-\ell_3^3\text{m}_1^3}{1-\ell_3\text{m}_1}\\[2ex]\dfrac{1-\ell_1^3\text{m}_2^3}{1-\ell_1\text{m}_2}&\dfrac{1-\ell_2^3\text{m}_2^3}{1-\ell_2\text{m}_2}&\dfrac{1-\ell_3^3\text{m}_2^3}{1-\ell_3\text{m}_2}\\[2ex]\dfrac{1-\ell_1^3\text{m}_3^3}{1-\ell_1\text{m}_3}&\dfrac{1-\ell_2^3\text{m}_3^3}{1-\ell_2\text{m}_3}&\dfrac{1-\ell_3^3\text{m}_3^3}{1-\ell_3\text{m}_3}\end{vmatrix}$}\par
here $\ell_1$, $\ell_2$, $\ell_3$, m$_1$, m$_2$, m$_3$ $\in$ R and $\ell_\text{i}$m$_\text{j}$ $\neq$ 1 $\forall$ i, j $\in$ \{1, 2, 3\} also $\ell_1<\ell_2<\ell_3$ \& m$_1<$ m$_2<$ m$_3$, then choose the INCORRECT option(s).
\begin{qoptions}
\opt{A}{D $>$ 0}
\opt{B}{D $<$ 0}
\opt{C}{D = 0}
\opt{D}{D = $-$1}
\end{qoptions}
\end{bankq}
\begin{bankq}[id=SA10-P1-MATH-S3-Q05, source={Score Advanced - Test 10 (06 May 2023)}, paper=1, qno=5, type=numerical, subtopic={Properties}, answer={52}, answer_src=key, difficulty=8, diff_parts={K3 S1 I2 C2 F1}, diff_note={concepts: maximum determinant with greatest integer entries, properties, bounds of the integer parts | steps: bounds of the integer parts}]
If 2 $\leq$ [x] $<$ 5, $-$2 $\leq$ [y] $<$ 4, 3 $\leq$ [z] $<$ 5, then maximum value of determinant $\left|\begin{matrix} [\text{x}]+2 & [\text{y}] & [\text{z}] \\ [\text{x}] & [\text{y}]+2 & [\text{z}] \\ [\text{x}] & [\text{y}] & [\text{z}]+2 \end{matrix}\right|$ is :\par
(where [.] denotes greatest integer function)
\end{bankq}
\begin{bankq}[id=SA11-P1-MATH-S3-Q03, source={Score Advanced - Test 11 (09 May 2023)}, paper=1, qno=3, type=numerical, subtopic={Properties}, answer={7}, answer_src=key, difficulty=8, diff_parts={K3 S1 I2 C2 F1}, diff_note={concepts: determinant with roots of a cubic, properties, row operations then symmetric functions | steps: row operations then symmetric functions}, also_in={Score Advanced - Test 29 (01 May 2024) P1 Q3}]
Let a, b, c are roots of x$^3$ + x$^2$ $-$ 2x + 1 = 0 then the value of $\left|\begin{matrix} a+c & b & a+2b+2c \\ a+b & c & b+2a+2c \\ b+c & a & c+2a+2b \end{matrix}\right|$ is equal to
\end{bankq}
\begin{bankq}[id=SA24-P1-MATH-S3-Q04, source={Score Advanced - Test 24 (17 Apr 2024)}, paper=1, qno=4, type=numerical, subtopic={Properties}, answer={1}, answer_src=key, difficulty=8, diff_parts={K3 S1 I2 C2 F1}, diff_note={concepts: value of a determinant at a given angle, properties, expand the determinant | steps: expand the determinant}]
Find the value of $f\left(\dfrac\pi6\right)$, where $f(\theta)=\left|\begin{array}{ccc}\cos^2\theta&\cos\theta\sin\theta&-\sin\theta\\ \cos\theta\sin\theta&\sin^2\theta&\cos\theta\\ \sin\theta&-\cos\theta&0\end{array}\right|$
\end{bankq}
\begin{bankq}[id=SA25-P2-MATH-S2-Q05, source={Score Advanced - Test 25 (21 Apr 2024)}, paper=2, qno=5, type=multi, subtopic={Properties}, answer={A,B,C,D}, answer_src=key, difficulty=8, diff_parts={K3 S1 I2 C2 F1}, diff_note={concepts: determinant condition from a matrix equation, properties, manipulate the equation with A and B | steps: manipulate the equation with A and B}, also_in={Score Advanced - Test 26 (21 Apr 2024) P2 Q5}]
Let A and B be two 3 $\times$ 3 matrices where both A and I + AB are non-singular. If $3A^2B^2+AB=2(AB)^2$ then, which of the following is greater than $|B|$ (where I is identity matrix of order 3 \& $|X|$ denotes determinant of X)
\begin{qoptions}
\opt{A}{1}
\opt{B}{2}
\opt{C}{3}
\opt{D}{4}
\end{qoptions}
\end{bankq}
\begin{bankq}[id=SA27-P2-MATH-S1-Q03, source={Score Advanced - Test 27 (24 Apr 2024)}, paper=2, qno=3, type=multi, subtopic={Properties}, answer={C}, answer_src=solved, working={Expanding the two determinants gives -x(x\^{}4 - 2x\^{}3 - 3x\^{}2 + 20x + 4) = 0; the quartic factor has two real and two non-real roots, matching option C (the printed options match the quartic; the extra factor x is possibly a transcription artefact).}, difficulty=8, diff_parts={K3 S1 I2 C2 F1}, diff_note={concepts: number of solutions of an equation of two determinants, properties, expand both determinants | steps: expand both determinants}]
The equation\par
\resizebox{\linewidth}{!}{$\begin{vmatrix}(1+x)^2&(1+x)^2&-(2+x^3)\\2x+1&3x&1-5x\\x+1&2x&2-3x\end{vmatrix}+\begin{vmatrix}(1+x)^2&2x+1&x+1\\(1-x)^2&3x&2x\\1-2x&3x-2&2x-3\end{vmatrix}=0$}\par
has
\begin{qoptions}
\opt{A}{No real solution}
\opt{B}{Exactly 4 real solutions}
\opt{C}{Exactly two real and two non-real solutions}
\opt{D}{Infinite number of solutions, real or non-real}
\end{qoptions}
\end{bankq}
\begin{bankq}[id=SA33-P1-MATH-S4-Q05, source={Score Advanced - Test 33 (12 May 2024)}, paper=1, qno=5, type=numerical, subtopic={Properties}, answer={8}, answer_src=key, difficulty=8, diff_parts={K3 S1 I2 C2 F1}, diff_note={concepts: standard problem in this chapter, properties, apply the relevant results | steps: apply the relevant results; compute}, also_in={Score Advanced - Test 37 (19 May 2024) P1 Q5}]
A 3 $\times$ 3 determinant has entries either 1 or $-$1. Let S be the set of all determinants such that the product of elements of any row or column is $-$1. For example $\begin{vmatrix}1&-1&1\\1&1&-1\\-1&1&1\end{vmatrix}$ is an element of S and number of elements of S is m. Let P = $\begin{vmatrix}3&-2&3\\2&-2&3\\0&-1&1\end{vmatrix}$ and trace of the matrix adj (adj P) is n then the value of $\dfrac mn$ is
\end{bankq}
\begin{bankpassage}[id=SA39-P1-MATH-S3-G1-2, source={Score Advanced - Test 39 (20 Apr 2025)}, questions={1-2}]
Let a, b and c be real numbers such that the system of linear equations\par
10x + 11y + 12z = a\par
13x + 14y + 15z = b\par
16x + 17y + 18z = c $-$ 3\par
is consistent. Let $|$x$|$ denotes the determinant of matrix\par
X = $\begin{bmatrix}\mathrm a&2&1\\\mathrm b&1&0\\\mathrm c&0&-1\end{bmatrix}$\par
Let P be the plane containing all those (a, b, c) for which system of linear equations is consistent and D be the square of the distance of the point ($-$9, 0, 0) from the plane P.
\end{bankpassage}
\begin{bankq}[id=SA39-P1-MATH-S3-Q01, source={Score Advanced - Test 39 (20 Apr 2025)}, paper=1, qno=1, type=numerical, subtopic={Properties}, answer={7}, answer_src=key, difficulty=8, diff_parts={K3 S1 I2 C2 F1}, diff_note={concepts: standard problem in this chapter, properties, apply the relevant results | steps: apply the relevant results; compute}, group=SA39-P1-MATH-S3-G1-2]
The value of ($|$x$|$ + 10) is \underline{\hspace{1.2cm}}
\end{bankq}
\begin{bankq}[id=SA39-P1-MATH-S3-Q02, source={Score Advanced - Test 39 (20 Apr 2025)}, paper=1, qno=2, type=numerical, subtopic={Properties}, answer={24}, answer_src=key, difficulty=8, diff_parts={K3 S1 I2 C2 F1}, diff_note={concepts: standard problem in this chapter, properties, apply the relevant results | steps: apply the relevant results; compute}, group=SA39-P1-MATH-S3-G1-2]
The value of D is \underline{\hspace{1.2cm}}
\end{bankq}

\section*{Enthusiast Pre-Score papers}
\begin{bankq}[id=EPS03-P1-MATH-S1-Q06, source={Enthusiast Pre-Score - Test 03 (23 Nov 2020)}, paper=1, qno=6, type=multi, subtopic={Properties}, answer={A,B,C,D}, answer_src=key, difficulty=7, diff_parts={K2 S2 I2 C1 F1}, diff_note={concepts: determinant of linear terms, properties | steps: row operations; simplify; evaluate sum and product}]
If $T_r=\left|\begin{array}{ccc}r-1&r-2&r-6\\2r-4&2r-6&2r-11\\3r-9&3r-12&3r-18\end{array}\right|$, then -
\begin{qoptions}
\opt{A}{$\displaystyle\sum_{r=0}^{5}T_r=36$}
\opt{B}{$\displaystyle\sum_{r=1}^{5}T_r=30$}
\opt{C}{$\displaystyle\prod_{r=0}^{5}T_r=6^6$}
\opt{D}{$\displaystyle\prod_{r=1}^{5}T_r=6^5$}
\end{qoptions}
\end{bankq}
\begin{bankq}[id=EPS03-P1-MATH-S3-Q06, source={Enthusiast Pre-Score - Test 03 (23 Nov 2020)}, paper=1, qno=6, type=numerical, subtopic={Expansion}, answer={2.00}, answer_src=key, difficulty=8, diff_parts={K2 S2 I2 C2 F1}, diff_note={concepts: determinant expansion matching a polynomial, expand the determinant | steps: expand the determinant; compare coefficients; solve}]
The value of `$\lambda$' if $ax^4+bx^3+cx^2+50x+d=\left|\begin{array}{ccc}x^3-14x^2&-x&3x+\lambda\\4x+1&3x&x-4\\-3&4&0\end{array}\right|$ is -
\end{bankq}
\begin{bankq}[id=EPS03-P2-MATH-S1-Q06, source={Enthusiast Pre-Score - Test 03 (23 Nov 2020)}, paper=2, qno=6, type=multi, subtopic={Properties}, answer={C}, answer_src=key, difficulty=7, diff_parts={K2 S2 I2 C1 F1}, diff_note={concepts: homogeneous system with parameters, properties | steps: determinant condition; factor; solutions}]
$\alpha x+y+z=0$, $x-\beta y+3z=0$ and $x-y-z=0$ has a non trivial solution, then
\begin{qoptions}
\opt{A}{$\alpha=1$ or $\beta=-5$}
\opt{B}{$\alpha=-1$ or $\beta=3$}
\opt{C}{$\alpha=-1$ or $\beta=-3$}
\opt{D}{$\alpha=1$ or $\beta=3$}
\end{qoptions}
\end{bankq}
\begin{bankq}[id=EPS03-P2-MATH-S3-Q01, source={Enthusiast Pre-Score - Test 03 (23 Nov 2020)}, paper=2, qno=1, type=numerical, subtopic={Expansion}, answer={5.00}, answer_src=key, difficulty=7, diff_parts={K2 S2 I2 C1 F1}, diff_note={concepts: maximum of a determinant, expansion | steps: expand; combine trig terms; maximum}]
If $\Delta=\begin{vmatrix}1&3\cos\theta&1\\ \sin\theta&1&3\cos\theta\\ 1&\sin\theta&1\end{vmatrix}$, then value of $\dfrac{\Delta_{\max}}{2}$ is.
\end{bankq}
\begin{bankq}[id=EPS03-P2-MATH-S3-Q05, source={Enthusiast Pre-Score - Test 03 (23 Nov 2020)}, paper=2, qno=5, type=numerical, subtopic={Expansion}, answer={6.00}, answer_src=key, difficulty=7, diff_parts={K2 S2 I2 C1 F1}, diff_note={concepts: roots of a determinant equation, expansion | steps: expand; polynomial; sum of roots}]
The sum of all the roots of equation $\begin{vmatrix}x(x-2)&1&2\\ 3&(x+2)(x-4)&4\\ 5&6&(x+4)(x-6)\end{vmatrix}=0$ will be-.
\end{bankq}

\section*{TOAS papers}
\begin{bankq}[id=TOAS07-P1-MATH-S1-Q02, source={TOAS - Test 07 (01 May 2022)}, paper=1, qno=2, type=single, subtopic={Properties}, answer={B}, answer_src=key, difficulty=4, diff_parts={K1 S2 I1 C1 F0}, diff_note={concepts: determinant coefficient of x | steps: f'(0) as a sum of determinants with one row differentiated; evaluate each; add | I: coefficient of x equals f'(0)}, also_in={TAAS - Test 01 (18 Jun 2023) P1 Q2; TNAS - Test 04 (18 Jun 2023) P1 Q2}]
If $f(x)=\begin{vmatrix}(x-1)^2&(x-2)^3&x+1\\(x-1)^3&x^2&(x+1)^2\\x+2&x+1&x^5\end{vmatrix}$ then coefficient of $x$ in $f(x)$ is -
\begin{qoptions}
\opt{A}{$-14$}
\opt{B}{$-16$}
\opt{C}{0}
\opt{D}{25}
\end{qoptions}
\end{bankq}
\begin{bankq}[id=TOAS07-P2-MATH-S4-Q08, source={TOAS - Test 07 (01 May 2022)}, paper=2, qno=8, type=numerical, subtopic={Properties}, answer={64}, answer_src=key, difficulty=5, diff_parts={K2 S2 I1 C0 F1}, diff_note={concepts: determinant identity, Vieta's relations | steps: sum of roots is 4; row operations give (alpha + beta + gamma)\^{}3; evaluate | I: determinant equals the cube of the sum}]
If $\alpha,\beta,\gamma$ are the roots of $4x^3-16x^2+19x-6=0$, then the value of the determinant
\[ \begin{vmatrix}\alpha-\beta-\gamma&2\alpha&2\alpha\\2\beta&\beta-\gamma-\alpha&2\beta\\2\gamma&2\gamma&\gamma-\alpha-\beta\end{vmatrix} \]
is
\end{bankq}
\begin{bankq}[id=TOAS15-P1-MATH-S3-Q02, source={TOAS - Test 15 (22 Aug 2022)}, paper=1, qno=2, type=numerical, subtopic={Properties}, answer={6.00}, answer_src=key, difficulty=3, diff_parts={K1 S2 I0 C0 F1}, diff_note={concepts: determinant of a product of scalar multiples | steps: determinant of a scalar multiple of a 2 by 2 matrix; powers; exponent}]
Consider the matrix $A,B,C,D$ with order $2\times3$, $3\times4$, $4\times4$, $4\times2$ respectively.\par
Let $x=(\alpha AB\gamma C^2D)^3$, where $\alpha$ \& $\gamma$ are scalars. Let $|x|=k|ABC^2D|^3$ \& if $k=(\alpha\gamma)^\lambda$, then $\lambda$ is (where $|P|$ denotes determinant value of matrix $P$)
\end{bankq}
\begin{bankq}[id=TOAS15-P1-MATH-S3-Q06, source={TOAS - Test 15 (22 Aug 2022)}, paper=1, qno=6, type=numerical, subtopic={Properties}, answer={0}, answer_src=solved, working={$|B|=3^{\sum i-\sum j}|A|=|A|$ and likewise $|C|=|B|$, so $|B|-|C|=0$}, difficulty=3, diff_parts={K1 S1 I1 C0 F1}, diff_note={concepts: determinant under row and column scaling | steps: factors cancel in the determinant; both determinants equal that of A}]
Let $A=[a_{ij}]_{3\times3}$, $B=[b_{ij}]_{3\times3}$ where $b_{ij}=3^{i-j}a_{ij}$, $C=[c_{ij}]_{3\times3}$, where $c_{ij}=4^{i-j}b_{ij}$ be any three matrices. If $|A|=4$, then $|B|-|C|$ ($|X|$ denotes determinant of matrix $X$)
\end{bankq}
\begin{bankq}[id=TOAS17-P2-MATH-S1-Q05, source={TOAS - Test 17 (18 Sep 2022)}, paper=2, qno=5, type=single, subtopic={Properties}, answer={A}, answer_src=key, difficulty=4, diff_parts={K1 S2 I1 C1 F0}, diff_note={concepts: determinant condition for non-trivial solution | steps: determinant equal to zero; solve for f(x); monotonicity}]
Consider the three linear equations with respect to $X$, $Y$ and $Z$ as
\par\noindent$x.X+Y+2Z=0$\par\noindent$f(x).X+3Y+x^2Z=0$\par\noindent$5xX+6y+Z=0$ having
\par\noindent non-trivial solution for $X$, $Y$ and $Z$ then curve $f(x)$
\begin{qoptions}
\opt{A}{is always increasing}
\opt{B}{is always decreasing}
\opt{C}{has exactly one critical point}
\opt{D}{has three critical points}
\end{qoptions}
\end{bankq}
\begin{bankq}[id=TOAS18-P1-MATH-S2-Q07, source={TOAS - Test 18 (30 Oct 2022)}, paper=1, qno=7, type=single, subtopic={Properties}, answer={C}, answer_src=solved, working={I: $BA$ gives $5x=1$, $x=1/5$ (Q); II: $(\alpha^2-4)^3=125$, $\alpha=3$ (P); III: $\Delta=a^2b^2c^2\times4$ (S); IV: $f(\theta)=1$ (R)}, difficulty=4, diff_parts={K1 S2 I1 C1 F0}, diff_note={concepts: determinants and inverse matrices | steps: find x from the product; alpha from the determinant of the cube; determinant with sum of squares zero; trigonometric determinant}]
\begin{center}\footnotesize\setlength{\tabcolsep}{1pt}
\begin{tabular}{|L{0.6cm}L{4.9cm}|L{0.6cm}L{0.7cm}|}\hline
&List-I&&List-II\\\hline
(I)&$A=\begin{bmatrix}2&0&7\\0&1&0\\1&-2&1\end{bmatrix}$ is inverse of $B=\begin{bmatrix}-x&14x&7x\\0&1&0\\1&-4x&-2x\end{bmatrix}$, then the value of $x$ is ?&(P)&3\\\hline
(II)&If $A=\begin{bmatrix}\alpha&2\\2&\alpha\end{bmatrix}$ and $|A^3|=125$ then $\alpha$ is equal to ($\alpha>0$)&(Q)&$\dfrac15$\\\hline
(III)&If $a,b,c$ are three complex numbers such that $a^2+b^2+c^2=0$ and $\Delta=\begin{vmatrix}b^2+c^2&ab&ac\\ab&c^2+a^2&bc\\ac&bc&a^2+b^2\end{vmatrix}=ka^2b^2c^2$&(R)&1\\\hline
(IV)&If $f(\theta)=\begin{vmatrix}\cos^2\theta&\cos\theta\sin\theta&-\sin\theta\\\cos\theta\sin\theta&\sin^2\theta&\cos\theta\\\sin\theta&-\cos\theta&0\end{vmatrix}$ then $f\left(\dfrac{2\pi}3\right)$ is&(S)&4\\\hline
&&(T)&2\\\hline
\end{tabular}\end{center}
\begin{qoptions}
\opt{A}{I $\to$ R;II $\to$ R;III $\to$ T;IV $\to$ T}
\opt{B}{I $\to$ R;II $\to$ P;III $\to$ T;IV $\to$ R}
\opt{C}{I $\to$ Q;II $\to$ P;III $\to$ S;IV $\to$ R}
\opt{D}{I $\to$ Q;II $\to$ P;III $\to$ S;IV $\to$ T}
\end{qoptions}
\end{bankq}

\section*{TAAS papers}
\begin{bankq}[id=TAAS01-P1-MATH-S3-Q01, source={TAAS - Test 01 (18 Jun 2023)}, paper=1, qno=1, type=numerical, subtopic={Properties}, answer={64.00}, answer_src=key, difficulty=6, diff_parts={K2 S2 I1 C1 F1}, diff_note={concepts: determinant of a structured matrix, roots of a cubic | steps: simplify the determinant using the sum of the roots; symmetric function values; evaluate}, also_in={TNAS - Test 04 (18 Jun 2023) P1 Q1}]
If $\alpha,\beta,\gamma$ are the roots of $4x^3-16x^2+19x-6=0$, then the value of the determinant $\begin{vmatrix}\alpha-\beta-\gamma&2\alpha&2\alpha\\2\beta&\beta-\gamma-\alpha&2\beta\\2\gamma&2\gamma&\gamma-\alpha-\beta\end{vmatrix}$ is
\end{bankq}

\section*{TOPS papers}
\begin{bankq}[id=TOPS03-P1-MATH-S1-Q02, source={TOPS - Test 03 (20 Nov 2022)}, paper=1, qno=2, type=single, subtopic={Properties}, answer={C}, answer_src=key, difficulty=5, diff_parts={K2 S2 I1 C1 F0}, diff_note={concepts: determinant, harmonic progression | steps: determinant zero condition; row reduction; relation among a b c}]
If the system of equations $x+4ay+az=0$, $x+3by+bz=0$ and $x+2cy+cz=0$ has a non-trivial solution, then $a,b,c$ are in ($a,b,c$ are distinct non-zero real numbers) -
\begin{qoptions}
\opt{A}{A.P.}
\opt{B}{G.P.}
\opt{C}{H.P.}
\opt{D}{A.G.P.}
\end{qoptions}
\end{bankq}
\begin{bankq}[id=TOPS03-P1-MATH-S2-Q07, source={TOPS - Test 03 (20 Nov 2022)}, paper=1, qno=7, type=multi, subtopic={Cramer's rule}, answer={A,C}, answer_src=key, difficulty=6, diff_parts={K2 S2 I1 C1 F1}, diff_note={concepts: determinant, inequality | steps: non zero solution means the determinant vanishes; relation among a b c; bound the expressions; check each}]
If the system of equations $x-cy-bz=0$, $cx-y+az=0$ and $bx+ay-z=0$ has a non-zero solution and atleast one of $a,b,c$ is a proper fraction, then
\begin{qoptions}
\opt{A}{$a^2+b^2+c^2<3$}
\opt{B}{$abc<-1$}
\opt{C}{$abc>-1$}
\opt{D}{$a^2+b^2+c^2>3$}
\end{qoptions}
\end{bankq}
\begin{bankpassage}[id=TOPS03-P1-MATH-S3-G1-2, source={TOPS - Test 03 (20 Nov 2022)}, questions={1-2}]
Consider a system of equations :

$x\sin\phi-2y\cos\phi-\lambda z=0$

$x+2y+z=0$

$-x+y+z=0$ (where $\lambda\in R$)
\end{bankpassage}
\begin{bankq}[id=TOPS03-P1-MATH-S3-Q01, source={TOPS - Test 03 (20 Nov 2022)}, paper=1, qno=1, type=numerical, subtopic={Cramer's rule}, answer={3}, answer_src=solved, working={The determinant gives $\sin\phi+4\cos\phi=3\lambda$, so $|3\lambda|\le\sqrt{17}$ and $|\lambda|\le1.37$; integers $-1,0,1$: three values}, difficulty=6, diff_parts={K2 S2 I1 C1 F1}, diff_note={concepts: determinant, range of a trigonometric function | steps: expand the determinant; range of sin plus cos; integer values}, group=TOPS03-P1-MATH-S3-G1-2]
The number of integer values of $\lambda$ for which system have non-trivial solution --
\end{bankq}
\begin{bankq}[id=TOPS03-P1-MATH-S3-Q02, source={TOPS - Test 03 (20 Nov 2022)}, paper=1, qno=2, type=numerical, subtopic={Cramer's rule}, answer={3}, answer_src=solved, working={$\lambda=1$ is attainable so statement 1 holds; $\lambda=3$ makes the determinant nonzero so the solution is unique (trivial) and statement 2 holds; three integral values so statement 3 holds. All three correct}, difficulty=5, diff_parts={K2 S1 I1 C1 F1}, diff_note={concepts: system of equations, uniqueness | steps: each statement with the determinant; count the correct ones}, group=TOPS03-P1-MATH-S3-G1-2]
Consider the following statements -

Statement 1 : For $\lambda=1$ there exist $\phi$ for which the system of equation will have infinite solutions.

Statement 2 : For $\lambda=3$ there exist $\phi$ for which the system of equation will have unique solution.

Statement 3 : Number of integral values of $\lambda$ is 3 for the system of equation to have non-trivial solutions.

Then, number of correct statements out of the above, is
\end{bankq}
\begin{bankq}[id=TOPS03-P1-MATH-S4-Q07, source={TOPS - Test 03 (20 Nov 2022)}, paper=1, qno=7, type=numerical, subtopic={Properties}, answer={1}, answer_src=key, difficulty=8, diff_parts={K2 S3 I2 C1 F1}, diff_note={concepts: determinant, trigonometric conditions | steps: row operations; the factored form; conditions on the angles; maximum alpha; evaluate}]
If
\par\noindent\hfill$\begin{vmatrix}1&1&1\\\cos\alpha&\cos(\alpha+\beta)&\cos(\alpha+\beta+\gamma)\\\sin\alpha&\sin(\alpha+\beta)&\sin(\alpha+\beta+\gamma)\end{vmatrix}=0$\hfill\null\par
such that $\alpha\in[0,2\pi]$ and $\beta,\gamma\in(0,2\pi)$, then $\dfrac{(\beta+\gamma)}{\alpha_{\max}}$ is equal to
\end{bankq}
\begin{bankq}[id=TOPS03-P2-MATH-S1-Q01, source={TOPS - Test 03 (20 Nov 2022)}, paper=2, qno=1, type=multi, subtopic={Properties}, answer={A,B}, answer_src=key, difficulty=7, diff_parts={K2 S2 I2 C1 F1}, diff_note={concepts: determinant properties, products | steps: row operations; value of T n; products and sums; check each}, also_in={TOPS - Test 04 (18 Dec 2022) P1 Q5; Star Review (Other) - Test 02 (10 Oct 2022) P1 Q4}]
If $T_n=\begin{vmatrix}n-1&n-2&n-6\\2n-4&2n-6&2n-11\\3n-9&3n-12&3n-18\end{vmatrix}$ where $n\in N$, then which of the following is/are true ?
\begin{qoptions}
\opt{A}{$\displaystyle\prod_{n=1}^{m}T_n=6^m$}
\opt{B}{$\displaystyle\sum_{n=1}^{10}T_n=60$}
\opt{C}{$\dfrac{T_{n+1}}{T_n}=2$}
\opt{D}{$T_{100}.T_{101}=48$}
\end{qoptions}
\end{bankq}
\begin{bankq}[id=TOPS03-P2-MATH-S3-Q03, source={TOPS - Test 03 (20 Nov 2022)}, paper=2, qno=3, type=numerical, subtopic={Properties}, answer={7}, answer_src=key, difficulty=9, diff_parts={K2 S3 I2 C2 F1}, diff_note={concepts: determinant, symmetric functions | steps: factor the determinant; use the relation between roots; sum of squares; evaluate the determinant; square root | cubic roots in a determinant}]
If $\alpha,\beta,\gamma$ are the roots of $x^3+x^2-2x+1=0$, then the value of $\sqrt\Delta$ is equal to \underline{\hspace{1.5cm}}

where $\Delta=\begin{vmatrix}\gamma^2&\beta^2&2\beta\gamma-\alpha^2\\2\alpha\gamma-\beta^2&\alpha^2&\gamma^2\\\alpha^2&2\alpha\beta-\gamma^2&\beta^2\end{vmatrix}$
\end{bankq}
\begin{bankpassage}[id=TOPS04-P1-MATH-S3-G5-6, source={TOPS - Test 04 (18 Dec 2022)}, questions={5-6}]
Consider the system of equations $2x+\lambda y+z=0$, $\lambda x+2y+z=0$ and $3x+4y+2z=0$.
\end{bankpassage}
\begin{bankq}[id=TOPS04-P1-MATH-S3-Q05, source={TOPS - Test 04 (18 Dec 2022)}, paper=1, qno=5, type=numerical, subtopic={Cramer's rule}, answer={-8.00}, answer_src=key, difficulty=8, diff_parts={K2 S3 I2 C1 F1}, diff_note={concepts: homogeneous system, determinant | steps: determinant zero; least lambda; solve the system; ratio; multiply}, group=TOPS04-P1-MATH-S3-G5-6]
For the least value of $\lambda$ if $(x_1,y_1,z_1)$ is a non trivial solution then $\dfrac{x_1+y_1}{z_1}$ is $K$ then value of $14K$ is equal to
\end{bankq}
\begin{bankq}[id=TOPS04-P1-MATH-S3-Q06, source={TOPS - Test 04 (18 Dec 2022)}, paper=1, qno=6, type=numerical, subtopic={Cramer's rule}, answer={49.00}, answer_src=key, difficulty=6, diff_parts={K2 S2 I1 C1 F1}, diff_note={concepts: determinant, roots | steps: determinant polynomial; sum of the roots; multiply}, group=TOPS04-P1-MATH-S3-G5-6]
The sum of all values of $\lambda$ for which the system has a non trivial solution is $P$ then value of $14P$ is equal to
\end{bankq}
\begin{bankq}[id=TOPS04-P2-MATH-S1-Q05, source={TOPS - Test 04 (18 Dec 2022)}, paper=2, qno=5, type=single, subtopic={Cramer's rule}, answer={C}, answer_src=key, difficulty=3, diff_parts={K1 S1 I1 C1 F0}, diff_note={concepts: homogeneous system | steps: determinant zero; solve m}]
The system of equation

$2x-y+z=0$

$x-2y+z=0$

$mx-y+2z=0$ has infinite number of non trivial solution for
\begin{qoptions}
\opt{A}{$m=2$}
\opt{B}{$m=-5$}
\opt{C}{$m=5$}
\opt{D}{$m=-2$}
\end{qoptions}
\end{bankq}
\begin{bankq}[id=TOPS04-P2-MATH-S2-Q02, source={TOPS - Test 04 (18 Dec 2022)}, paper=2, qno=2, type=numerical, subtopic={Expansion}, answer={3}, answer_src=solved, working={The determinant is $-2x^2+4x-3$, product of roots $3/2$, so $2k=3$}, difficulty=6, diff_parts={K2 S2 I1 C1 F1}, diff_note={concepts: determinant, Vieta | steps: expand; quadratic in x; product}]
If product of the roots of $\begin{vmatrix}1+x&1&1\\3&1&2\\3&x&3\end{vmatrix}=0$ is $k$, then the value of $2k$ is
\end{bankq}
\begin{bankq}[id=TOPS05-P1-MATH-S3-Q01, source={TOPS - Test 05 (14 Jan 2023)}, paper=1, qno=1, type=numerical, subtopic={Cramer's rule}, answer={42.00}, answer_src=key, difficulty=6, diff_parts={K2 S2 I1 C1 F1}, diff_note={concepts: determinant, count | steps: determinant value; unique condition; count pairs}]
Given $a,b\in\{0,1,2,\ldots,6\}$. Consider the system of equations

$x+y+z=4$,

$2x+y+3z=6$,

$x+2y+az=b$

The number of ordered pairs $(a,b)$ so that the system of equations has unique solution, is
\end{bankq}

\section*{TMAS papers}
\begin{bankq}[id=TMAS04-P2-MATH-S2-Q02, source={TMAS - Test 04 (19 Nov 2023)}, paper=2, qno=2, type=single, subtopic={Properties}, answer={2}, answer_src=key, difficulty=5, diff_parts={K1 S2 I2 C1 F0}, diff_note={concepts: determinant of a circulant type matrix | steps: structure of the coefficient matrix; determinant; scaling}]
If $a,b,c$ be the distinct real numbers and $x,y,z$ be any three variables such that the system of linear equations

$(3z-2y)a+(2x-z)b+(y-3x)c=0$

$(y-3x)a+(3z-2y)b+(2x-z)c=0$

$(2x-z)a+(y-3x)b+(3z-2y)c=0$

can be express as $AX=O$, $A$ is $3\times3$ square matrix and $X=\begin{bmatrix}x\\y\\z\end{bmatrix}$ then the value of $\det(A)+\det(2A)+\det(3A)$ is :
\end{bankq}
\begin{bankq}[id=TMAS04-P2-MATH-S3-Q06, source={TMAS - Test 04 (19 Nov 2023)}, paper=2, qno=6, type=numerical, subtopic={Properties}, answer={9.85}, answer_src=key, difficulty=6, diff_parts={K1 S2 I2 C1 F1}, diff_note={concepts: counting singular matrices | steps: determinant zero condition; products equal; count patterns}]
Let $M=[a_{ij}]_{2\times2}$ whose all entries are distinct and $a_{ij}\in\{1,2,5,10\}$, then number of such singular matrices $M$ is
\end{bankq}

\section*{TNAS papers}
\begin{bankq}[id=TNAS07-P1-MATH-S2-Q09, source={TNAS - Test 07 (20 Aug 2023)}, paper=1, qno=9, type=numerical, subtopic={Properties}, answer={2}, answer_src=key, difficulty=6, diff_parts={K1 S2 I2 C1 F1}, diff_note={concepts: determinant of a transformed matrix | steps: row sums structure; matrix relation; determinant; value}]
Let $P=[p_{ij}]$ be a $3\times3$ non singular matrix and a matrix $Q=[q_{ij}]$ is formed such that $q_{ij}=\displaystyle\sum_{\substack{r=1\\r\neq j}}^{3}p_{ir}$ then value of $\det(P^{-1}Q)$ is
\end{bankq}

\section*{Star Review (Other) papers}
\begin{bankq}[id=STAR02-P1-MATH-S3-Q03, source={Star Review (Other) - Test 02 (10 Oct 2022)}, paper=1, qno=3, type=numerical, subtopic={Properties}, answer={1.50}, answer_src=key, difficulty=8, diff_parts={K2 S2 I2 C2 F1}, diff_note={concepts: determinant with powers, count of integer triplets | steps: factor the determinant; simplify; solve for the sum of powers; count triplets}]
If number of possible triplet $(x,y,z)$, which satisfies the given determinant, where $x,y,z\in I^+$\par
$\begin{vmatrix}x^3+1&x^2y&x^2z\\xy^2&y^3+1&y^2z\\xz^2&yz^2&z^3+1\end{vmatrix}=11$ is $k$ then value of $\dfrac k2$ is
\end{bankq}
\begin{bankq}[id=STAR02-P1-MATH-S3-Q07, source={Star Review (Other) - Test 02 (10 Oct 2022)}, paper=1, qno=7, type=numerical, subtopic={Expansion}, answer={9.00}, answer_src=key, difficulty=7, diff_parts={K1 S2 I2 C2 F1}, diff_note={concepts: determinant of a function and its maximum | steps: expand; simplify to a quadratic; maximum value using the given difference}]
Let $a_2,a_3\in R$ such that $|a_2-a_3|=6$ and $f(x)=\begin{vmatrix}1&a_3&a_2\\1&a_3&2a_2-x\\1&2a_3-x&a_2\end{vmatrix}$, $x\in R$. Then the greatest value of $f(x)$ is
\end{bankq}
\begin{bankq}[id=STAR02-P2-MATH-S2-Q06, source={Star Review (Other) - Test 02 (10 Oct 2022)}, paper=2, qno=6, type=multi, subtopic={Properties}, answer={D}, answer_src=key, difficulty=7, diff_parts={K1 S2 I2 C2 F1}, diff_note={concepts: determinant of a symmetric expression | steps: expand; factor; compare with the given form; value of k}]
If $\begin{vmatrix}2bc-a^2&c^2&b^2\\c^2&2ca-b^2&a^2\\b^2&a^2&2ab-c^2\end{vmatrix}=(a^3+b^3+c^3+kabc)^2$, then $k$ is equal to
\begin{qoptions}
\opt{A}{$2$}
\opt{B}{$-2$}
\opt{C}{$3$}
\opt{D}{$-3$}
\end{qoptions}
\end{bankq}
\begin{bankq}[id=STAR02-P2-MATH-S2-Q10, source={Star Review (Other) - Test 02 (10 Oct 2022)}, paper=2, qno=10, type=multi, subtopic={Expansion}, answer={A,D}, answer_src=key, difficulty=5, diff_parts={K1 S2 I1 C1 F1}, diff_note={concepts: determinant equation | steps: expand; factor; roots}]
If $\begin{vmatrix}1&4&20\\1&-2&5\\1&2x&5x^2\end{vmatrix}=0$, then values of $x$ satisfying the equation is/are
\begin{qoptions}
\opt{A}{$2$}
\opt{B}{$-2$}
\opt{C}{$1$}
\opt{D}{$-1$}
\end{qoptions}
\end{bankq}
\begin{bankq}[id=STAR02-P2-MATH-S3-Q01, source={Star Review (Other) - Test 02 (10 Oct 2022)}, paper=2, qno=1, type=numerical, subtopic={Properties}, answer={2}, answer_src=key, difficulty=7, diff_parts={K1 S2 I2 C2 F1}, diff_note={concepts: number of real solutions of a determinant equation | steps: factor powers of x; reduce to a polynomial; solve; count}]
The total number of distinct $x\in R$ for which $\begin{vmatrix}x&x^2&1+x^3\\2x&4x^2&1+8x^3\\3x&9x^2&1+27x^3\end{vmatrix}=10$ is
\end{bankq}
\begin{bankq}[id=STAR02-P2-MATH-S3-Q02, source={Star Review (Other) - Test 02 (10 Oct 2022)}, paper=2, qno=2, type=numerical, subtopic={Properties}, answer={4}, answer_src=key, difficulty=7, diff_parts={K1 S2 I2 C2 F1}, diff_note={concepts: maximum of a determinant in sine | steps: expand; quadratic in sine; maximum}]
Maximum value of determinant {\footnotesize$\begin{vmatrix}1&\sin\theta&1\\-\sin\theta&1&\sin\theta\\-1&-\sin\theta&1\end{vmatrix}$} is
\end{bankq}
\end{document}
